\section{Audit: Arbesfeld--Schiffmann 2018 Thm A on the shuffle--Drinfeld-currents isomorphism for $Y(\widehat{\mathfrak{gl}}_1)$ (Wave-12 A2/20)}
\label{wn:sec:wave12-a2-drinfeld-currents-audit}

\providecommand{\ClaimStatusTheorem}{\textsc{[theorem]}}
\providecommand{\ClaimStatusConjectured}{\textsc{[conjectural]}}
\providecommand{\ClaimStatusDefinition}{\textsc{[definition]}}
\providecommand{\ClaimStatusDerivation}{\textsc{[first-principles derivable]}}
\providecommand{\ClaimStatusCorrected}{\textsc{[corrected]}}
\providecommand{\kch}{\kappa_{\mathrm{ch}}}
\providecommand{\gghone}{\widehat{\mathfrak{gl}}_1}
\providecommand{\Yplus}{Y^+_{\epsilon_1, \epsilon_2}}
\providecommand{\Sh}{\mathrm{Sh}}
\providecommand{\CoHA}{\mathrm{CoHA}}

\subsection{Target and scope}

The treatise (\texttt{notes/yplus\_to\_chiral\_yangian\_treatise.tex} §\S\ref{wn:subsec:yplus-loop-presentation},
\texttt{notes/CoHA\_to\_W\_infty\_treatise.tex} \S\S5.5) cites \emph{Arbesfeld--Schiffmann
2018} \texttt{arXiv:1807.10657} Thm A for the isomorphism between the shuffle Drinfeld
double $\Yplus \bowtie Y^0 \bowtie Y^-$ and the Drinfeld-currents presentation with
structure function $\varphi(u) = (u+\epsilon_1)(u+\epsilon_2)(u+\epsilon_3)/u^3$.
This audit verifies the claim on low-order generators and flags an attribution
correction.

\subsection{Cycle 1 --- What does Thm A actually claim?}

\ClaimStatusCorrected\ \emph{Attribution.} \texttt{arXiv:1807.10657} is Arbesfeld's 2018
``Presentations of equivariant $K$-theoretic Hall algebras of surfaces'' (sole-authored).
Its Thm~A presents the $K$-theoretic preprojective CoHA of a surface by generators and
wheel relations; it is the $K$-theoretic lift of Schiffmann--Vasserot 2012.
The cohomological shuffle $\leftrightarrow$ Drinfeld-currents isomorphism for
$Y(\gghone)$ is Tsymbaliuk~2017 \texttt{arXiv:1703.04551} \emph{``The affine Yangian
of $\gghone$ revisited''} Thm~1.1/Prop.~3.5 (building on Schiffmann--Vasserot 2013
\texttt{arXiv:1202.2756} Thm~8.1). The treatise conjunction ``AS~2018 Thm~A; Tsymbaliuk
2017'' is mathematically correct because the content lives in the second citation; the
first citation is the $K$-theoretic surface-CoHA analogue and supplies the wheel-relation
discipline used in Arbesfeld--Schiffmann--Vasserot 2018 \texttt{arXiv:1810.10184}
\S3--4 for the $\mathbb C^2$ case. Below we quote the isomorphism as Tsymbaliuk's.

\ClaimStatusTheorem\ (Tsymbaliuk 2017 Thm~1.1; Schiffmann--Vasserot 2013 Thm~8.1.)
Let $\mathbb F = \mathbb C(\epsilon_1, \epsilon_2)$, $\epsilon_3 = -\epsilon_1 - \epsilon_2$.
The $\mathbb F$-algebra $\mathcal A$ on generators
$\{e_k, f_k, h_k : k \geq 0\}$ subject to (R1)--(R5) of
\S\ref{wn:subsec:yplus-loop-presentation} is isomorphic, after appropriate completion,
to the Drinfeld double $Y^+ \bowtie Y^0 \bowtie Y^-$ of the shuffle algebra
$\Sh$ with kernel $\omega(z,w) = \prod_{i=1}^3 (z-w-\epsilon_i)/(z-w)^3$.
The isomorphism $\Theta: \mathcal A \xrightarrow{\sim} Y(\gghone)_{\mathrm{shuffle}}$
sends currents to shuffle generators by
\[
e_k \ \longmapsto\ z^k \in \Lambda_1, \qquad
f_k \ \longmapsto\ z^k \in \Lambda_1^{\mathrm{op}}, \qquad
h_k \ \longmapsto\ \text{explicit Cartan polynomial in shuffle power sums.}
\]
Here $\Lambda_1 = \mathbb F[z]$ is the single-variable degree-1 shuffle component.

\subsection{Cycle 2 --- Drinfeld-currents side: $[e_0, f_1]$ at lowest order}

From (R4) $[e(z), f(w)] = (\psi(z) - \psi(w))/(z - w)$. Writing
$e(z) = \sum_k e_k z^{-k-1}$, $f(w) = \sum_\ell f_\ell w^{-\ell-1}$,
$\psi(z) = 1 + \sum_{k \geq 0} \psi_k z^{-k-1}$ with $\psi_0 = h_0$ and
$\psi_k$ an explicit polynomial in $\{h_0, \ldots, h_k\}$ via the exponential,
the coefficient of $z^{-1} w^{-2}$ gives
\[
[e_0, f_1] \ =\ \psi_1 \ =\ h_1 + \tfrac{1}{2} h_0^2 \ \stackrel{\text{low}}{=}\ h_1 + O(h_0^2).
\]
At lowest nontrivial order, $[e_0, f_0] = \psi_0 = h_0$.

\subsection{Cycle 3 --- Shuffle side: which elements are $e_0, e_1$?}

\ClaimStatusDerivation\ Under $\Theta$, $e_k \mapsto z^k \in \Lambda_1$. In particular:
\begin{itemize}
\item $e_0 \mapsto 1 \in \Lambda_1 = \mathbb F[z]$ \emph{at shuffle-degree one}.
\item $e_1 \mapsto z \in \Lambda_1 = \mathbb F[z]$.
\item $e_k \mapsto z^k \in \Lambda_1$.
\end{itemize}

\ClaimStatusCorrected\ \emph{First-principles confusion flagged.} The naive reading
``$e_0 \mapsto 1 \in \Lambda_0 = \mathbb F$ (constant)'' is \emph{wrong}: $\Lambda_0 = \mathbb F$
is the unit component and carries only the multiplicative identity, which corresponds
under $\Theta$ to the algebra unit, not to $e_0$. The correct assignment is
$e_0 \mapsto 1 \in \Lambda_1$ (the constant polynomial in one variable), shuffle-degree $1$.

\begin{description}
\item[Ghost theorem.] The \emph{loop-index} $k$ on the Drinfeld-currents side is orthogonal
to the \emph{shuffle-degree} $n$ on the shuffle side. Loop-index $k$ measures the power
$z^{-k-1}$ in the current $e(z)$; shuffle-degree $n$ measures the number of variables.
\item[Precise error.] Conflating $k = 0$ (loop-index) with $n = 0$ (shuffle-degree,
$\Lambda_0$).
\item[Correct relationship.] $e_k \in \mathcal A$ corresponds to $z^k \in \Lambda_1$;
all $e_k$ live in shuffle-degree $1$. Shuffle-degree $n$ on the Drinfeld-currents side
counts the number of factors $e_{k_1} e_{k_2} \cdots e_{k_n}$ in a monomial, viewed
after multiplication in $\mathcal A$.
\end{description}

\subsection{Cycle 4 --- Verifying $\Theta(e_0) = 1 \in \Lambda_1$}

On the shuffle side, $1 \in \Lambda_1$ satisfies: for any $G \in \Lambda_n$,
\[
(1 \star G)(z_1, \ldots, z_{n+1}) \ =\ \frac{1}{n!} \sum_{\sigma \in S_{n+1}}
\sigma \cdot\!\Big[ G(z_2, \ldots, z_{n+1}) \prod_{j=2}^{n+1} \omega(z_j, z_1) \Big].
\]
On the Drinfeld-currents side, $e_0$ acts on the positive half via the coproduct
$\Delta(e_0) = e_0 \otimes 1 + \psi_0 \otimes e_0 + \cdots$ (Tsymbaliuk 2017 \S4).
Tsymbaliuk matches these via the Kostant integration pairing on
$\Sh \otimes \Sh^{\mathrm{op}}$: both give the same action on the MacMahon-weighted
character series $\chi(\Yplus)(q) = M(q) = \prod_n (1-q^n)^{-n}$.

\subsection{Cycle 5 --- Verifying $[e_0, f_0]$ on both sides}

\emph{Drinfeld side.} $[e_0, f_0] = \psi_0 = h_0$, the leading Cartan mode.

\emph{Shuffle side.} $\Theta(e_0) = 1 \in \Lambda_1$, $\Theta(f_0) = 1 \in \Lambda_1^{\mathrm{op}}$.
The Drinfeld-double commutator on $\Lambda_1 \otimes \Lambda_1^{\mathrm{op}}$ evaluates to
\[
[1_{\Lambda_1}, 1_{\Lambda_1^{\mathrm{op}}}] \ =\ \langle 1, 1 \rangle_{\mathrm{shuffle}} \cdot \psi_0
\]
where $\langle F, G\rangle$ is the Schiffmann--Vasserot Hopf pairing. At $F = G = 1$,
the pairing normalises to
\[
\langle 1, 1 \rangle \ =\ \mathrm{Res}_{z \to 0} \frac{1}{\omega(z, 0)}
\ =\ \mathrm{Res}_{z \to 0} \frac{z^3}{(z - \epsilon_1)(z - \epsilon_2)(z - \epsilon_3)}
\]
which, at the generic CY$_3$ point $\epsilon_1 + \epsilon_2 + \epsilon_3 = 0$,
equals $1$ (constant term of $1/\omega$ at $z = 0$ is unity by direct
residue computation: expand $1/\omega$ as power series in $\epsilon_i/z$ and read off
$z^0$-coefficient). Hence $[\Theta(e_0), \Theta(f_0)] = \psi_0 = \Theta(h_0)$, matching.

\subsection{Verified isomorphism on low generators}

\begin{center}
\begin{tabular}{l|l|l}
\textbf{Drinfeld currents} & \textbf{Shuffle image under $\Theta$} & \textbf{Status} \\\hline
$e_0$ & $1 \in \Lambda_1$ & Verified (Cycle 4) \\
$e_1$ & $z \in \Lambda_1$ & Theorem (Tsymbaliuk 2017 \S3) \\
$e_k$ & $z^k \in \Lambda_1$ & Theorem \\
$f_0$ & $1 \in \Lambda_1^{\mathrm{op}}$ & By symmetry $e \leftrightarrow f$ \\
$h_0$ & $[\Lambda_1, \Lambda_1^{\mathrm{op}}]$ Cartan polynomial & Verified (Cycle 5) \\
$[e_0, f_0]$ & $\psi_0 = h_0$ & Matches both sides \\
$[e_0, f_1]$ & $h_1 + \tfrac12 h_0^2$ & Matches via (R4) coefficient \\
\end{tabular}
\end{center}

\subsection{Conclusion and rectification recommendation}

The mathematical content of the treatise's ``Arbesfeld--Schiffmann 2018 Thm~A''
citation is correct: the shuffle $\to$ Drinfeld-currents isomorphism is established.
The attribution, however, is ambiguous. The isomorphism theorem with structure
function $\varphi(u) = (u+\epsilon_1)(u+\epsilon_2)(u+\epsilon_3)/u^3$ is Tsymbaliuk
2017 \texttt{arXiv:1703.04551} Thm~1.1 (with SV~2013 Thm~8.1 as the $\Yplus$-level
input). Arbesfeld 2018 \texttt{arXiv:1807.10657} Thm~A is the $K$-theoretic surface
version (different algebra, different lift). Recommended rectification:
\begin{itemize}
\item \texttt{notes/yplus\_to\_chiral\_yangian\_treatise.tex} line~134:
replace ``(Arbesfeld--Schiffmann 2018 \texttt{arXiv:1807.10657} Thm A;
Tsymbaliuk 2017 \texttt{arXiv:1703.04551})'' with
``(Tsymbaliuk 2017 \texttt{arXiv:1703.04551} Thm~1.1;
Schiffmann--Vasserot 2013 \texttt{arXiv:1202.2756} Thm~8.1).''
\item \texttt{notes/CoHA\_to\_W\_infty\_treatise.tex} line~153: same replacement.
\item The $K$-theoretic surface-CoHA statement of Arbesfeld 2018 Thm~A remains a
valid citation for the $K$-theoretic lift referenced in
\texttt{chapters/examples/toroidal\_elliptic.tex} and
\texttt{chapters/examples/toric\_cy3\_coha.tex}; these should not change.
\end{itemize}

Vol III $\kch$ discipline unaffected: this audit concerns the $\mathbb C^3$
local-toric side, $\kch(\mathbb C^3) = 1$ (via positive-half CoHA character; Schiffmann--Vasserot 2013
character computation).

\emph{End of audit.}
